% ============================================================================
% CHAPTER 8 TEACHER'S MANUAL
% Complete instructional guide for teaching HITL system architecture
% ============================================================================
% Source: /markdown/ch8/Ch8T_updated.md
% Note: This file is for the Instructor's Edition, not the main book
% ============================================================================

\chapter{Teacher's Manual: Chapter 8}
\label{ch:teacher-ch08}

\begin{chapterquote}{Complete instructional guide}
Teaching HITL system architecture using \emph{The Blueprint for Smart Helper Systems}.
\end{chapterquote}

% ============================================================================
\section{Course Integration Guide}
% ============================================================================

\subsection{Learning Objectives}

By the end of this chapter, students will be able to:

\paragraph{Knowledge (Remembering \& Understanding).}
\begin{itemize}
    \item Identify the seven components of a complete HITL system architecture
    \item Distinguish between HITL pipelines (no learning) and HITL loops (learning)
    \item Describe four feedback integration approaches and their trade-offs
    \item Name six common production failure modes and their architectural causes
    \item Describe the HITL system lifecycle from prototype to monitoring-dominant operation
\end{itemize}

\paragraph{Application \& Analysis.}
\begin{itemize}
    \item Apply the seven-component framework to analyze a described system
    \item Diagnose which component is absent or deficient in a described failure
    \item Compare batch retraining, online learning, active learning, and RLHF for a given use case
    \item Explain why removing any single component breaks something important
    \item Analyze the Uber Tempe incident as an architectural failure
\end{itemize}

\paragraph{Synthesis \& Evaluation.}
\begin{itemize}
    \item Design a complete seven-component HITL architecture for a new domain
    \item Propose monitoring metrics for each architectural component
    \item Critique a described HITL system for architectural gaps
    \item Evaluate the trade-offs between different feedback integration strategies
\end{itemize}

\subsection{Prerequisites}
\begin{itemize}
    \item \textbf{Chapters 1--7}: Five-dimension framework, uncertainty quantification, interface design
    \item \textbf{Chapter 4}: Threshold band design (the routing layer builds on this)
    \item \textbf{Chapter 6}: Interface design (the interface layer references this in depth)
    \item \textbf{Optional}: Basic software architecture concepts helpful for the component discussion
\end{itemize}

\subsection{Course Positioning}
\begin{description}
    \item[Systems Course (Week 8--10)] Core architectural synthesis chapter; teach after all component concepts are established
    \item[Applied ML Course] Excellent bridge from theoretical concepts to production system design
    \item[Advanced Course] Focus on failure mode taxonomy and monitoring strategy; use design exercise
\end{description}

% ============================================================================
\section{Key Concepts for Instruction}
% ============================================================================

\subsection{The Pipeline vs.\ Loop Distinction}

This is the chapter's foundational concept and must be established clearly before the component taxonomy.

\begin{table}[htbp]
\caption{HITL Pipeline vs.\ HITL Loop: The critical distinction}
\label{tab:pipeline-vs-loop}
\centering
\begin{tabular}{@{}p{3.5cm}p{4cm}p{4cm}@{}}
\toprule
\textbf{Aspect} & \textbf{Pipeline} & \textbf{Loop} \\
\midrule
Human output goes\ldots & Downstream (to user/decision) & Back into model training \\
Model improves? & Never & Over time \\
Required extra components & None & Feedback integration, monitoring, audit \\
Failure mode & Permanent errors & Drift, gaming, error propagation \\
Example & Email triage (static model) & Spam filter that learns \\
\bottomrule
\end{tabular}
\end{table}

\begin{technote}[Teaching Tip]
Draw this on the board. Students consistently confuse ``a human reviews outputs'' with ``the system learns from humans.'' The diagram makes the distinction tangible. A pipeline has one arrow direction. A loop has arrows going both ways.
\end{technote}

\subsection{The Seven-Component Architecture}

\begin{table}[htbp]
\caption{The seven components of HITL architecture}
\label{tab:seven-components}
\centering
\begin{tabular}{@{}p{0.5cm}p{3cm}p{4cm}p{4cm}@{}}
\toprule
\textbf{\#} & \textbf{Component} & \textbf{Role} & \textbf{What breaks without it} \\
\midrule
1 & Prediction & Model output + confidence & Nothing to route; no system \\
2 & Uncertainty Quantification & Calibrated confidence scores & Routing sends wrong cases to humans \\
3 & Routing & Three-zone decision logic & Humans overwhelmed or system fully automated \\
4 & Interface & Human review UX & Human feedback is garbage in, garbage out \\
5 & Feedback Integration & Loop closure & System never improves (pipeline not loop) \\
6 & Monitoring & Ongoing behavior tracking & Drift is invisible; degradation uncaught \\
7 & Audit & Decision logging & System cannot be debugged or held accountable \\
\bottomrule
\end{tabular}
\end{table}

\subsection{The Four Feedback Integration Strategies}

\begin{table}[htbp]
\caption{Feedback integration strategies compared}
\label{tab:feedback-strategies}
\centering
\begin{tabular}{@{}p{2.5cm}p{2cm}p{2cm}p{2cm}p{3cm}@{}}
\toprule
\textbf{Strategy} & \textbf{Latency} & \textbf{Complexity} & \textbf{Error Risk} & \textbf{Best For} \\
\midrule
Batch retraining & Slow (monthly) & Low & Low & Stable, well-defined tasks \\
Online learning & Fast (real-time) & Medium & High & Rapidly evolving tasks \\
Active learning & Varies & High & Medium & Scarce labels, evolving distribution \\
RLHF & Training-time & Very high & Medium & LLMs, preference-based tasks \\
\bottomrule
\end{tabular}
\end{table}

% ============================================================================
\section{Lecture Planning Guide}
% ============================================================================

\subsection{50-Minute Lecture Structure}

\subsubsection{Opening (8 minutes): The Uber Tempe Incident}

\paragraph{How to present this case.}
\begin{itemize}
    \item Present clearly: the system detected the pedestrian; uncertainty was quantified; classification cycled through multiple wrong categories; the system eventually settled on a confident incorrect classification
    \item The backup safety driver was present but was looking at a video stream; no uncertainty signal was communicated to her in actionable form
    \item Key argument: calling something ``human-in-the-loop'' doesn't make it so. The loop requires a working \emph{connection} between uncertainty detection, human notification, and human capacity to act. Any broken link fails the whole loop
\end{itemize}

\begin{keyconcept}[Central Argument]
The Uber Tempe incident is not a sensor failure or a human failure. It is an \emph{architecture} failure. Both the algorithm and the human were present. Neither was connected to the other in a way that enabled effective intervention.
\end{keyconcept}

\subsubsection{Part 1 (10 minutes): Pipeline vs.\ Loop}

\paragraph{Board exercise.}
Draw the pipeline, then draw the loop. Ask:
\begin{itemize}
    \item ``What happens to a pipeline that processes incorrect cases? It keeps being incorrect.''
    \item ``What happens to a loop?'' (It gets better --- or worse, if the feedback is noisy)
    \item ``What's required to close the loop?'' (Components 5, 6, 7)
\end{itemize}

\subsubsection{Part 2 (14 minutes): Seven Components}

Walk through each component:
\begin{enumerate}
    \item \textbf{Prediction}: model output + confidence (covered deeply in earlier chapters)
    \item \textbf{Uncertainty Quantification}: calibrated scores, temperature scaling
    \item \textbf{Routing}: the three-zone system from Chapter 4, implemented here
    \item \textbf{Interface}: Chapter 6's domain; role in the architecture is receiving and structuring human feedback
    \item \textbf{Feedback Integration}: the loop-closing layer (will expand in next segment)
    \item \textbf{Monitoring}: keeps the system honest over time
    \item \textbf{Audit}: accountability, debugging, legal compliance
\end{enumerate}

\begin{trythis}
For each component, ask: ``Which component of the Uber Tempe system was absent or deficient?'' Students should identify Component 3 (routing / escalation) as the primary failure, with Component 4 (interface) as a contributing failure.
\end{trythis}

\subsubsection{Part 3 (10 minutes): Feedback Integration Options}

Use a concrete example for each strategy:
\begin{itemize}
    \item \textbf{Batch retraining}: spam filter updated monthly
    \item \textbf{Online learning}: live product recommendation system
    \item \textbf{Active learning}: medical annotation workflow with rare label categories
    \item \textbf{RLHF}: language model fine-tuning using preference comparisons
\end{itemize}

\subsubsection{Part 4 (10 minutes): Production Failure Modes}

Work through the failure taxonomy:
\begin{description}
    \item[Reviewer bottleneck] Queue grows faster than it is reviewed; high-stakes cases age; SLA violations
    \item[Label quality degradation] Human reviewers tired or disengaged; noise enters training data
    \item[Feedback lag] Model slow to incorporate corrections; human compensates for errors that model doesn't know it's making
    \item[Distribution shift without detection] World changes; model doesn't; no monitoring catches it
    \item[Ghost reviewer] Review role exists on paper; no one is actually doing it
    \item[Specification gaming] Model optimizes the feedback metric, not the underlying task
\end{description}

\begin{trythis}
For each failure mode: ``Which component of the seven would detect this? Which prevents it?'' Students should practice mapping failures to components.
\end{trythis}

\subsubsection{Close (8 minutes)}
Summary and assign architecture design exercise.

\subsection{75-Minute Extended Session}

\subsubsection{Deep Dive: Monitoring (15 minutes)}
\begin{itemize}
    \item Prediction monitoring: track distribution of confidence scores over time
    \item Human response monitoring: agreement rates, time-per-review, override rates
    \item Calibration drift: are confidence scores still calibrated?
    \item Distribution shift: MMD (Maximum Mean Discrepancy) and KL divergence approaches
    \item When monitoring should trigger: threshold exceedance vs.\ trend analysis
\end{itemize}

\subsubsection{Design Workshop (20 minutes)}
Students work in groups to design the seven-component architecture for one of:
\begin{itemize}
    \item A fraud detection system at a mid-sized bank
    \item A content moderation system for a news comment section
    \item A medical imaging triage system
    \item An email routing system for a legal department
\end{itemize}

% ============================================================================
\section{Discussion Questions by Level}
% ============================================================================

\subsection{Introductory Level}

\begin{enumerate}
    \item \textbf{Pipeline vs.\ Loop}: ``What is the key difference between a HITL pipeline and a HITL loop? Give one real-world example of each.''

    \item \textbf{Architecture Failure}: ``Why does the Uber Tempe incident qualify as an architecture failure rather than a sensor failure or a driver failure? Which architectural component was missing or deficient?''

    \item \textbf{Monitoring Role}: ``Describe two failure modes that Component 6 (Monitoring) would detect. Why can't the other components detect these failure modes instead?''
\end{enumerate}

\subsection{Intermediate Level}

\begin{enumerate}
    \setcounter{enumi}{3}
    \item \textbf{New Threat Scenario}: ``A content moderation system uses monthly batch retraining. A new type of harmful content emerges in week 1. Describe the system's behavior from week 1 to week 5. What architecture change would reduce this lag?''

    \item \textbf{Feedback Strategy Comparison}: ``Compare active learning and batch retraining on: latency, label efficiency, implementation complexity, and error propagation risk. For what specific use case would you choose each?''

    \item \textbf{Override Rate Analysis}: ``A HITL system's audit logs show human reviewers have been overriding automated decisions at 40\% over the past month. Describe three hypotheses for this. How would you diagnose which one is correct?''
\end{enumerate}

\subsection{Advanced Level}

\begin{enumerate}
    \setcounter{enumi}{6}
    \item \textbf{Complete Architecture Design}: ``Design the complete seven-component architecture for a system detecting potentially problematic emails in a legal department inbox. Specify each component, the feedback integration strategy, monitoring metrics, and escalation protocols. Identify the three highest-risk failure modes.''

    \item \textbf{Specification Gaming}: ``Specification gaming occurs when the model learns to optimize the feedback metric rather than the underlying task. Describe a concrete scenario, identify which architectural component is responsible for catching it, and propose a design intervention.''
\end{enumerate}

% ============================================================================
\section{Hands-On Activities \& Exercises}
% ============================================================================

\subsection{Activity 1: Component Identification Hunt (15 minutes)}

\textbf{Setup:} Students work individually or in pairs with descriptions of three real systems.

\textbf{Task:} For each system, identify which of the seven components are present and which are absent.

\paragraph{Systems to analyze.}
\begin{itemize}
    \item A bank's fraud alert system that notifies customers by SMS but never uses responses to update its model
    \item An email spam filter that retrains weekly on items users manually mark as spam
    \item A customer service chatbot that routes complex queries to human agents but doesn't record agent responses
\end{itemize}

\paragraph{Debrief questions.}
\begin{itemize}
    \item Which component is most commonly missing?
    \item What failure mode does each missing component create?
\end{itemize}

\subsection{Activity 2: Feedback Strategy Selection (20 minutes)}

\textbf{Setup:} Small groups receive a domain description.

\textbf{Task:} Select the appropriate feedback integration strategy and justify the choice using the four comparison dimensions (latency, label efficiency, complexity, error propagation risk).

\paragraph{Domain descriptions.}
\begin{description}
    \item[Medical imaging] Radiologists review uncertain scans; labels are expensive; distribution changes slowly; errors are high-stakes
    \item[Social media moderation] 10,000 posts/hour; new abuse patterns emerge daily; moderators work in shifts; labels are fast but noisy
    \item[Legal document review] Lawyers review flagged clauses; distribution is very stable; labels are authoritative; latency is acceptable
    \item[Recommendation system] User preferences are the feedback; personalization improves with each interaction; billions of data points daily
\end{description}

\subsection{Activity 3: Architecture Design Challenge (30 minutes)}

\textbf{Setup:} Groups of 3--4 students.

\textbf{Task:} Design the full seven-component architecture for a fraud detection system at a mid-sized bank.

\paragraph{Deliverables.}
\begin{itemize}
    \item Completed seven-component description (one paragraph per component)
    \item One key metric per component
    \item Top three predicted failure modes
    \item Feedback integration strategy with justification
\end{itemize}

\begin{table}[htbp]
\caption{Architecture design worksheet}
\label{tab:architecture-worksheet}
\centering
\begin{tabular}{@{}p{3.5cm}p{3cm}p{3cm}p{2.5cm}@{}}
\toprule
\textbf{Component} & \textbf{Design Specification} & \textbf{Key Metric} & \textbf{Failure Risk} \\
\midrule
Prediction & & & \\[6pt]
Uncertainty Quantification & & & \\[6pt]
Routing & & & \\[6pt]
Interface & & & \\[6pt]
Feedback Integration & & & \\[6pt]
Monitoring & & & \\[6pt]
Audit & & & \\
\bottomrule
\end{tabular}
\end{table}

% ============================================================================
\section{Assessment Strategies}
% ============================================================================

\subsection{Formative Assessment}

\paragraph{Exit Ticket.}
``Name the seven components of a HITL architecture and identify which one is most commonly missing in systems you've encountered.''

\paragraph{Quick Poll.}
``Which component closes the loop?'' (Answer: Feedback Integration)

\paragraph{Think-Pair-Share.}
``The Uber Tempe system had a human in the seat. Why wasn't she effectively in the loop?''

\subsection{Summative Assessment Options}

\subsubsection{Option 1: Architecture Analysis (500--750 words)}

\textbf{Prompt:} ``Examine the Uber Tempe incident post-mortem. Identify which of the seven architectural components was absent or deficient. Propose three specific architectural changes that could have prevented the incident, explaining which component each change addresses.''

\paragraph{Rubric.}
\begin{itemize}
    \item \textbf{Component Identification (35\%)}: Correctly identifies the deficient components with evidence from the case
    \item \textbf{Architectural Reasoning (30\%)}: Shows understanding of why the component failure caused the incident
    \item \textbf{Improvement Proposals (35\%)}: Proposes feasible, specific changes tied to named components
\end{itemize}

\subsubsection{Option 2: Full Architecture Design}

\textbf{Prompt:} ``Design the complete seven-component HITL architecture for a fraud detection system at a bank. For each component: (a) describe your design, (b) identify one key metric, and (c) identify the most likely production failure mode.''

\paragraph{Grade Bands.}
\begin{description}
    \item[A] All seven components present and well-specified; metrics are measurable; failure modes are specific and plausible; feedback strategy is justified
    \item[B] All seven components present; some underspecified; metrics mostly appropriate; failure modes identified
    \item[C] Five or six components present; some significant gaps; limited failure mode analysis
    \item[D] Fewer than five components; fundamental misunderstanding of pipeline vs.\ loop
\end{description}

% ============================================================================
\section{Common Student Misconceptions \& How to Address Them}
% ============================================================================

\subsection{Misconception 1: ``Having a human review step means the system is a HITL loop''}

\paragraph{Correction strategy.}
\begin{itemize}
    \item This is the pipeline mistake. The human review step must be connected to model improvement.
    \item Exercise: ask students whether their bank's fraud alert system --- where they confirm or deny a flag --- actually improves the model. Most don't know. That's the point.
    \item Framework connection: ``Feedback Integration is the component that distinguishes a loop from a pipeline.''
\end{itemize}

\subsection{Misconception 2: ``Active learning is just uncertain-case routing''}

\paragraph{Correction strategy.}
\begin{itemize}
    \item Active learning optimizes the \emph{learning value} of routed cases, not just their uncertainty
    \item A system that routes only uncertain cases doesn't explore the training distribution efficiently; it may over-sample a specific failure region
    \item Key distinction: routing by uncertainty vs.\ routing by information gain
\end{itemize}

\subsection{Misconception 3: ``Monitoring should catch errors in real time''}

\paragraph{Correction strategy.}
\begin{itemize}
    \item Some failure modes require real-time monitoring (reviewer bottleneck, high review rate)
    \item Others require trend analysis over weeks or months (slow drift, specification gaming)
    \item Different monitoring timescales are needed for different failure modes
\end{itemize}

\subsection{Misconception 4: ``Audit logs are just for compliance''}

\paragraph{Correction strategy.}
\begin{itemize}
    \item Audit trails are the primary tool for debugging production failures
    \item Without Component 7, you cannot determine what caused a bad decision
    \item Show: the Q6 scenario (40\% override rate) requires audit logs to diagnose
\end{itemize}

% ============================================================================
\section{Extension Activities \& Research Projects}
% ============================================================================

\subsection{Individual Research Topics}
\begin{enumerate}
    \item \textbf{Post-Mortem Analysis}: Research a publicized AI system failure (facial recognition, content moderation, medical AI). Apply the seven-component framework. Which component was deficient?
    \item \textbf{Production Architecture}: Research how a major tech company (Google, Amazon, Meta) has described its HITL architecture for a specific product. What components are documented?
    \item \textbf{RLHF Deep Dive}: Research the architecture of RLHF as used in ChatGPT or similar systems. Map it to the seven-component framework.
\end{enumerate}

\subsection{Group Projects}
\begin{enumerate}
    \item \textbf{Architecture Audit}: Select a deployed AI system, interview users or read public documentation, and assess which architectural components are present or absent.
    \item \textbf{Failure Mode Catalog}: Create a database of documented HITL failures mapped to the six failure mode taxonomy.
    \item \textbf{Monitoring Design}: Design a complete monitoring stack (Component 6) for a specific production system, including metrics, thresholds, alerting, and escalation procedures.
\end{enumerate}

% ============================================================================
\section{Connections to Other Chapters}
% ============================================================================

\subsection{Backward Links}
\begin{itemize}
    \item \textbf{Chapter 1}: Five-dimension framework is the analytical lens for all seven components
    \item \textbf{Chapter 4}: Threshold band design is implemented in Component 3 (Routing)
    \item \textbf{Chapter 6}: Interface design is Component 4; this chapter situates it architecturally
    \item \textbf{Chapter 7}: RLHF is Component 5 (one feedback integration strategy)
\end{itemize}

\subsection{Forward Links}
\begin{itemize}
    \item \textbf{Chapter 9}: Tools for implementing annotation (Component 4 UX) and feedback pipelines
    \item \textbf{Chapter 12}: Measurement framework provides the metrics for Component 6 (Monitoring)
    \item \textbf{Chapter 15}: Failure diagnosis chapter connects back to the failure mode taxonomy
    \item \textbf{Chapter 17}: Active learning in depth (Component 5 strategy)
\end{itemize}

\bigskip
\begin{center}
\rule{0.5\textwidth}{0.4pt}
\end{center}
\bigskip

\noindent\textit{Chapter~8 provides the architectural synthesis point for everything covered in Parts I--II. The seven-component framework is the scaffold students will carry into all subsequent application chapters. Invest time in the pipeline vs.\ loop distinction and the failure mode taxonomy --- these concepts recur throughout the book.}
